\documentclass[11pt,a4paper,oneside]{report}
\usepackage[T1]{fontenc}
\usepackage[utf8]{inputenc}
\usepackage[margin=26mm,headheight=14pt]{geometry}
\usepackage{amsmath,amssymb,amsthm,mathtools,bm,booktabs,array,longtable,microtype,xcolor}
\usepackage{enumitem}
\usepackage[colorlinks=true,linkcolor=blue!50!black,urlcolor=blue!55!black]{hyperref}
\usepackage{fancyhdr}
\pagestyle{fancy}\fancyhf{}\lhead{Quantum Gravity Dynamics}\rhead{Chapter 3: Normalization}\cfoot{\thepage}
\numberwithin{equation}{chapter}
\newtheorem{theorem}{Theorem}[chapter]
\newtheorem{proposition}[theorem]{Proposition}
\newtheorem{lemma}[theorem]{Lemma}
\newtheorem{corollary}[theorem]{Corollary}
\theoremstyle{definition}\newtheorem{definition}[theorem]{Definition}
\newtheorem{assumption}[theorem]{Assumption}
\theoremstyle{remark}\newtheorem{remark}[theorem]{Remark}
\newcommand{\dd}{\mathrm d}
\newcommand{\ii}{\mathrm i}
\newcommand{\hbarc}{\hbar c}
\newcommand{\Phist}{\Phi_{\mathrm{stat}}}
\newcommand{\Phin}{\Phi_{\mathrm N}}
\newcommand{\GNewton}{G_{\mathrm N}}
\newcommand{\epsQ}{\epsilon_{\mathrm Q}}
\setcounter{chapter}{2}
\title{\textbf{Quantum Gravity Dynamics}\\[5pt]\large Chapter 3 --- Gravitational Normalization, the Macroscopic Potential, and Quantum-Force Verification\\[8pt]\normalsize Proof-oriented research manuscript, QGD--095 coordinated edition}
\author{Research working draft}
\date{5 November 2026 --- Action-first normalization review (QGD--095 retained)}
\begin{document}
\maketitle
\tableofcontents
\clearpage
\chapter{Gravitational Normalization and Quantum-Force Verification}
\label{ch:normalization}
\begin{abstract}
This chapter starts with QGD's complete local macroscopic four-mode Dirac construction and its proposed signed natural-invariant master metric. In the explicitly restricted, classical, fixed-frame, spherical sector with an \emph{adopted} Einstein--Hilbert action, we derive the vacuum natural-invariant profile and fix its one physical exterior integration constant by Newtonian correspondence. We then derive the exact static Killing-energy potential per unit test mass, its all-order Taylor expansion, and an explicit remainder bound. The dimensionless geometric inverse-coupling convention involving a Compton length is distinguished from the original Dirac wavefunction multiplier. We give the corrected flux bridge, a solvable nonrelativistic radial wave, and a conditional curved-space radial Dirac expansion, retaining every explicit occurrence of Planck's constant. A separate referee-style audit shows that the historic universal attractive $9/2$ quantum-force coefficient has \emph{not} been derived: quantum corrections to a separated radial-wave equation are not automatically corrections to gravity. A new Dirac-to-Pauli quantum-source result fixes the density-gradient coefficient $\hbar^2/(8m_f)$ and derives a conditional quantum-matter-sourced gravitational response; Gauss's law forbids a universal exterior $r^{-4}$ correction from compact spherical source structure alone. Claims are classified by assumptions and proof obligations; exact symbolic verification does not constitute Lean-kernel certification.
\end{abstract}

\section*{Action-defined normalization: interpretation for the revised QGD paradigm}
\addcontentsline{toc}{section}{Action-defined normalization}
\noindent The independent gravitational variables of QGD are the admissible natural-invariant one-forms, with physical geometry fixed by Axiom~I. Their classical dynamics follows from the chosen Einstein--Hilbert--matter action and an ordinary stationary-action principle. Consequently $G_N$ and $\Lambda$ are calibrated physical couplings and boundary data, not quantities QGD must derive from microscopic quantum matter. This chapter's static solutions, Newtonian limit and Killing-energy potential are exact \emph{effective-gravity correspondence tests}. Historical speculative force coefficients are not additional axioms or universal QGD predictions. The independent Axiom~II of microscopic realizability is not used in the derivations in this chapter.


\section{Scope, conventions, and axioms}\label{sec:scope}
We use metric signature $(-+++)$, $x^0=ct$, $\hbar=h/(2\pi)>0$, $c>0$, $\GNewton>0$, an exterior source mass $M>0$, and, where needed, an independently specified positive test- or reference-fermion mass $m_t,m_{\rm ref}>0$. The Newtonian gravitational constant is an experimental correspondence input, not a number claimed to follow from microscopic Dirac dynamics. The dimensionless local Dirac ratios are denoted $\zeta_i$; the \emph{promoted gravitational natural invariants} retain QGD's original notation $\sigma_{\mu}^{(s)}$. These are distinct objects.

\begin{assumption}[Microscopic and macroscopic Dirac hand-off]\label{ass:dirac}
Chapter 1 supplies a specified microscopic Dirac effective action and, under stated coarse-graining and local/WKB hypotheses, a macroscopic Dirac-equivalent state. A local positive-energy momentum eigenmode satisfies
\begin{equation}
\bm\zeta=\frac{c\bm\pi}{E_{\rm loc}+m_{\rm eff}c^2},\qquad
\mathsf R=\bm\tau\cdot\bm\zeta,
\label{eq:ratio}
\end{equation}
where $\bm\tau$ are Pauli matrices. No assertion is made that one global outgoing mode exactly solves the full interacting, curved-space, macroscopic Dirac equation.
\end{assumption}

For a consistent common local momentum convention the complete four-mode state is
\begin{equation}
\boxed{\displaystyle
\Psi(x)=\frac{\alpha_G}{r}\left[
(\psi_0u_{\uparrow}+\psi_1u_{\downarrow})e^{-\ii S_+/\hbar}
+(\psi_2v_{\uparrow}+\psi_3v_{\downarrow})e^{+\ii S_-/\hbar}
\right],}
\label{eq:fourmode}
\end{equation}
with local columns
\begin{equation}
\begin{aligned}
 u_{\uparrow}&=\begin{pmatrix}1\\0\\\zeta_z\\\zeta_x+\ii\zeta_y\end{pmatrix},&
 u_{\downarrow}&=\begin{pmatrix}0\\1\\\zeta_x-\ii\zeta_y\\-\zeta_z\end{pmatrix},\\
 v_{\uparrow}&=\begin{pmatrix}-\zeta_z\\-\zeta_x-\ii\zeta_y\\1\\0\end{pmatrix},&
 v_{\downarrow}&=\begin{pmatrix}-\zeta_x+\ii\zeta_y\\\zeta_z\\0\\1\end{pmatrix}.
\end{aligned}
\label{eq:columns}
\end{equation}
The overall prefactor $\alpha_G$ is not observable independently of a compensating common rescaling of the four reference mode amplitudes. The radial factor $r^{-1}$ belongs to the stated ansatz, not to a universal normalization theorem.

\begin{assumption}[The complete QGD master-metric proposal]\label{ass:metric}
The proposed metric is
\begin{equation}
\boxed{\displaystyle
 g_{\mu\nu}=T^a{}_{\mu}T^b{}_{\nu}
 \left[\eta_{ab}-\sum_{s}\varepsilon_s\,
 \sigma_a^{(s)}\sigma_b^{(s)}-\lambda\mathcal Q_{ab}\right],
 \qquad \varepsilon_s\in\{-1,+1\}.}
\label{eq:master}
\end{equation}
The microscopic promotion from Dirac-derived observables to independent $\sigma_{\mu}^{(s)}$ fields is QGD's defining postulate. The precise derivative operator $\mathcal Q_{ab}$ and its coefficient $\lambda$ have not been microscopically matched or independently varied in the general model. They are not assumed to vanish as a theorem.
\end{assumption}

\begin{assumption}[Classical benchmark only]\label{ass:benchmark}
For the calculations in Sections \ref{sec:static}--\ref{sec:series} choose the explicit $\lambda=0$, fixed-asymptotically-flat-frame, one-positive-signed-sector, static spherical restriction of \eqref{eq:master}; use the \emph{adopted} Einstein--Hilbert action
\begin{equation}
 S_0=\frac{c^3}{16\pi\GNewton}\int\dd^4x\sqrt{-g}\,R
 +S_D[\Psi,g],
 \label{eq:action}
\end{equation}
where $\dd^4x=\dd x^0\,\dd^3x$ and $x^0=ct$. In the vacuum exterior $T_{\mu\nu}=0$. Variations are performed independently in the declared sector, with well-posed boundary conditions and the requisite gravitational boundary term. No microscopic derivation of the Einstein--Hilbert coefficient is claimed.
\end{assumption}

\begin{remark}[Action units]
To avoid a hidden convention change in variations below, only the nonzero overall multiplier of the Einstein--Hilbert functional matters. The radial reduced Lagrangian is therefore written with $\propto$. With $x^0=ct$, the customary dimensionful coefficient is $c^3/(16\pi G)$ for $\dd x^0\,\dd^3x$; equivalently it is $c^4/(16\pi G)$ for $\dd t\,\dd^3x$, since $\dd x^0=c\,\dd t$. Both give the same action and Euler--Lagrange equations. Physical $c$ factors are restored explicitly in the calibration and observables.
\end{remark}

\section{Exact classical spherical reduction}\label{sec:static}
Let $r$ be areal radius, and let $A,B$ be real $C^2$ functions on an unbounded connected exterior interval. With the designated positive signed sector, set
\begin{equation}
 \sigma=A(r)\dd x^0+B(r)\dd r,\quad
 F:=1-A^2,\quad D:=1+B^2-A^2.
 \label{eq:AB}
\end{equation}
Assume $F>0$, $D>0$, and $AB\ne0$ on this interval. Before coordinate diagonalization, the induced metric reads
\begin{equation}
 \dd s^2=-F(\dd x^0)^2+2AB\dd x^0\dd r
 +(1+B^2)\dd r^2+r^2\dd\Omega^2.
 \label{eq:crossmetric}
\end{equation}
A stationary coordinate change removes the cross term locally, yielding
\begin{equation}
 \dd s^2=-F\dd\widetilde x^{0\,2}+\frac DF\dd r^2+r^2\dd\Omega^2.
 \label{eq:diagmetric}
\end{equation}
The curvature calculation uses this equivalent metric; the independent $A,B$ variations refer to the declared spherical sector \emph{prior to imposing the solution}.

\begin{lemma}[Four-dimensional curvature and boundary identity]\label{lem:EH}
For \eqref{eq:diagmetric}, with primes denoting $\dd/\dd r$, the scalar curvature is
\begin{equation}
 R=-\frac{F''}{D}+\frac{F'D'}{2D^2}
 -\frac{4F'}{rD}+\frac{2FD'}{rD^2}
 +\frac2{r^2}-\frac{2F}{r^2D}.
 \label{eq:R}
\end{equation}
Furthermore, defining $L_{\rm red}:=(rF)'/\sqrt D+\sqrt D$, the exact identity is
\begin{equation}
 r^2\sqrt D\,R=2L_{\rm red}
 +\frac\dd{\dd r}\left(-\frac{r^2F'}{\sqrt D}-\frac{4rF}{\sqrt D}\right).
 \label{eq:EHidentity}
\end{equation}
\end{lemma}
\begin{proof}
Compute the Levi--Civita connection of the four-dimensional diagonal metric, contract to $R$, and differentiate the displayed boundary expression using the product rule. This proves an equality of functions for $r>0,F>0,D>0$ before any equations of motion. A complete tensor-level Lean proof remains outstanding.
\end{proof}

\begin{proposition}[Independent natural-invariant vacuum equations]\label{prop:EL}
Under Assumption \ref{ass:benchmark}, variation of $L_{\rm red}$ with respect to $A$ and $B$ independently gives
\begin{equation}
 \mathcal E_A=-\frac{A(rB^2)'}{D^{3/2}}=0,
 \qquad
 \mathcal E_B=\frac{B[B^2+r(A^2)']}{D^{3/2}}=0.
 \label{eq:EL}
\end{equation}
Consequently, on the nondegenerate exterior branch,
\begin{equation}
 \boxed{A^2(r)=B^2(r)=\frac{C_g}{r}},\qquad C_g>0.
 \label{eq:profile}
\end{equation}
\end{proposition}
\begin{proof}
The first Euler--Lagrange equation, with $A\ne0$, implies $(rB^2)'=0$. Thus $B^2=C_g/r$. The second, with $B\ne0$, gives $(A^2)'=-C_g/r^2$, hence $A^2=C_g/r+C_0$. Asymptotic flatness with $A^2\to0$ at infinity fixes $C_0=0$. Direct substitution gives the vacuum Schwarzschild metric $D=1$, $F=1-C_g/r$. \emph{Neither $A=B$ nor the numerical value of $C_g$ was assumed in taking the variations.}
\end{proof}

\begin{remark}[Exact scope of Proposition \ref{prop:EL}]
This is a restricted theorem of an adopted classical action. It does not prove that the single positive sector solves an arbitrary full signed-sector/tetrad variation with unmatched $\lambda\mathcal Q_{ab}$. The sector, smoothness, asymptotic normalization and boundary treatment are substantive hypotheses.
\end{remark}

\section{One Newtonian calibration and three distinct potentials}\label{sec:calibration}
Since $g_{00}=-1+C_g/r$, define the weak-field metric potential by
\begin{equation}
 \Phi_{\rm metric}(r):=\frac{c^2}{2}(g_{00}+1)
 =-\frac{c^2C_g}{2r}.
 \label{eq:metricpotential}
\end{equation}
The slow-particle weak-field radial acceleration is $-\Phi_{\rm metric}'$. Matching once to Newton's exterior acceleration $-\GNewton M/r^2$ fixes
\begin{equation}
 \boxed{C_g=\frac{2\GNewton M}{c^2}=:r_s.}
 \label{eq:calibration}
\end{equation}
The source mass $M$ is an independently defined gravitational mass, not an arbitrary test mass. The adopted Einstein--Hilbert coefficient already involves $G$; this correspondence identifies the physical normalization of that coefficient and the exterior source integration constant, not a prediction of $G$ from Dirac theory.

\begin{definition}[Static Killing-energy potential]\label{def:phi}
For a test mass $m_t>0$ held static at radius $r>r_s$, with timelike Killing vector normalized at infinity, let
\begin{equation}
 E_\infty(r)=m_tc^2\sqrt{-g_{00}(r)}
 =m_tc^2\sqrt{1-\frac{r_s}{r}},\qquad
 \boxed{\Phist(r):=\frac{E_\infty(r)-m_tc^2}{m_t}.}
 \label{eq:PhiDef}
\end{equation}
\end{definition}
Note that $\Phi_{\rm metric}=-GM/r$ for this coordinate representation, and $\Phin=-GM/r$ in the Newtonian exterior, but $\Phist$ differs beyond leading order. It is a supported-body energy per unit test mass; it is not a general position-only Hamiltonian for orbiting probes and it is not a quantum-corrected gravitational field.

\begin{theorem}[Exact calibrated static potential]\label{thm:phi}
For $r>r_s$, the static Killing-energy potential is
\begin{equation}
 \boxed{\displaystyle
 \Phist(r)=c^2\left[\sqrt{1-\frac{2\GNewton M}{c^2r}}-1\right]
 =-\frac{c^2(r_s/r)}{1+\sqrt{1-r_s/r}}.}
 \label{eq:PhiExact}
\end{equation}
\end{theorem}
\begin{proof}
Substitute the calibrated $g_{00}$ into Definition \ref{def:phi}. The second expression follows by multiplying numerator and denominator by $1+\sqrt{1-r_s/r}$. The leading term is $-GM/r$ by continuity as $r_s/r\to0$.
\end{proof}

\section{An all-order expansion, and a rigorously bounded remainder}\label{sec:series}
Define $x=r_s/r$ and $a_n=\binom{2n}{n}/[4^n(2n-1)]$ for $n\ge1$.
\begin{theorem}[Convergent exterior expansion]\label{thm:series}
For $0<x<1$ one has
\begin{equation}
 \boxed{\displaystyle\Phist(r)=-c^2\sum_{n=1}^\infty a_nx^n},
 \qquad \frac{a_{n+1}}{a_n}=\frac{2n-1}{2n+2}<1.
 \label{eq:allorder}
\end{equation}
In physical units,
\begin{equation}
 \boxed{\begin{aligned}
 \Phist(r)={}&-\frac{\GNewton M}{r}
 -\frac{\GNewton^2M^2}{2c^2r^2}
 -\frac{\GNewton^3M^3}{2c^4r^3}\\
 &-\frac{5\GNewton^4M^4}{8c^6r^4}+\mathcal R_4(r),
 \end{aligned}}
 \label{eq:fourterms}
\end{equation}
where
\begin{equation}
 \boxed{\displaystyle |\mathcal R_4(r)|\leq
 \frac{7c^2}{256}\frac{x^5}{1-x}.}
 \label{eq:remainder}
\end{equation}
\end{theorem}
\begin{proof}
Use the generalized binomial series for $\sqrt{1-x}$. Its nonconstant coefficients are $-a_n$ for $|x|<1$. Because $a_{n+1}/a_n<1$ and $a_5=7/256$, the magnitude of the tail after $n=4$ is bounded by $c^2(7/256)\sum_{n=5}^\infty x^n$, yielding \eqref{eq:remainder}. Insert $r_s=2GM/c^2$ in the first four coefficients for \eqref{eq:fourterms}.
\end{proof}

\begin{corollary}[Static radial force derived from the energy]\label{cor:force}
At fixed test mass $m_t$, the outward-positive force defined by the radial derivative of the static Killing energy is
\begin{equation}
 \boxed{\displaystyle
 F_{r,\mathrm{stat}}=-m_t\Phist'(r)
 =-\frac{\GNewton M m_t}{r^2\sqrt{1-r_s/r}}.}
 \label{eq:staticforce}
\end{equation}
This definition is the derivative of the conserved energy at infinity
with respect to the areal coordinate $r$.  In the Schwarzschild static
orthonormal frame, the proper holding-force magnitude is also
$m_tG_{\rm N}M/(r^2\sqrt F)$; their equality in this chart follows from
the compensating proper-distance and Killing-redshift factors.  One must
still specify the operational observer before calling a radial derivative
a measured force.  No new gravitational interaction is introduced by this
energy-gradient definition.
Its first terms are
\begin{equation}
 F_{r,\mathrm{stat}}=-\frac{\GNewton M m_t}{r^2}
 \left[1+\frac{\GNewton M}{c^2r}
 +\frac{3\GNewton^2M^2}{2c^4r^2}+\cdots\right].
\end{equation}
These are classical, not $\hbar$-dependent.
\end{corollary}

\section{The geometric normalization and the original $\alpha_G$}\label{sec:alpha}
The Dirac theory supplies a Compton length once a \emph{reference mass} $m_{\rm ref}>0$ has been independently selected:
\begin{equation}
 \bar\lambda_{\rm ref}=\frac{\hbar}{m_{\rm ref}c}.
\end{equation}
\begin{definition}[Dimensionless geometric inverse-strength convention]
\begin{equation}
 \boxed{\displaystyle
 |\alpha_G^{\rm geom}|^2
 :=\frac{\bar\lambda_{\rm ref}}{C_g}
 =\frac{\hbar c}{2\GNewton M m_{\rm ref}}.}
 \label{eq:geomalpha}
\end{equation}
\end{definition}
This is a precise \emph{definition after Newtonian calibration}, not a theorem relating the gravitational natural invariant to the common multiplier of the four-mode wavefunction. In particular,
\begin{equation}
 \boxed{\displaystyle
 \frac{\hbar}{m_{\rm ref}c|\alpha_G^{\rm geom}|^2}=C_g
 =\frac{2\GNewton M}{c^2}.}
 \label{eq:hcancellation}
\end{equation}
The $\hbar$ in this parametrization cancels identically from the classical potential \eqref{eq:PhiExact} and all terms of \eqref{eq:fourterms}. Its appearance in a chosen coordinate on parameter space cannot be counted as a quantum correction.

To see the normalization obstruction directly, let $\alpha_G$ multiply all the original Dirac amplitudes in \eqref{eq:fourmode}. For any nonzero complex $z$,
\begin{equation}
 \alpha_G\longmapsto z\alpha_G,\qquad
 \psi_j\longmapsto z^{-1}\psi_j\quad(j=0,1,2,3)
 \label{eq:rescaling}
\end{equation}
leaves the physical $\Psi$ unchanged. Its local component ratios remain unchanged as well. Hence the spinor prefactor is unidentifiable without an independently imposed reference-state norm or new physical coupling.

For a specified conserved occupation $N$ and chosen norm $N_{\rm ref}$, a linear homogeneous Dirac source gives the state-amplitude convention
\begin{equation}
 |\alpha_G^{\rm state}|^2=\frac{N}{N_{\rm ref}}.
 \label{eq:statealpha}
\end{equation}
The equality $\alpha_G^{\rm state}=\alpha_G^{\rm geom}$ is \emph{not} a conclusion of Chapters 1--2. For a model in which both magnitudes are equated by extra matching, the necessary compatibility condition is
\begin{equation}
 N_{\rm ref}=\frac{2\GNewton M m_{\rm ref}N}{\hbar c}.
 \label{eq:compatibility}
\end{equation}
It is a constraint that a future microscopic matching theory would need to justify; it is not a prediction of the present chapter.

\section{The corrected Chapter 1 radial current bridge}\label{sec:bridge}
For a coherent, outgoing, locally on-shell, positive-energy radial WKB branch, do not drop the lower Dirac components before evaluating its flux. Write
\begin{equation}
 \Psi\simeq\frac{\alpha_G\psi_0(r)}{r}e^{-\ii S/\hbar}
 \begin{pmatrix}\chi\\[2pt]\zeta(\bm\tau\cdot\widehat{\bm r})\chi\end{pmatrix},
 \quad \zeta=\frac{cp}{\Delta},\quad
 \Delta=E_{\rm loc}+m_{\rm eff}c^2,\quad\chi^\dagger\chi=1.
 \label{eq:WKB}
\end{equation}
Direct Dirac matrix algebra gives
\begin{equation}
 j^r=\frac{2c^2}{r^2}|\alpha_G\psi_0|^2\frac{p}{\Delta}.
 \label{eq:current}
\end{equation}
For stationary spherical outward flux $\dot N=4\pi r^2j^r>0$ in the declared local-flat/reference flux convention, this yields
\begin{equation}
 \boxed{\displaystyle
 |\alpha_G\psi_0(r)|^2=\frac{\dot N\,\Delta(r)}{8\pi c^2p(r)}
 =\mathcal C_J^2\Delta(r)\lambda_{\rm dB}(r)},
 \quad\mathcal C_J^2:=\frac{\dot N}{8\pi c^2h},\ \lambda_{\rm dB}:=\frac hp.
 \label{eq:bridge}
\end{equation}
The bridge is a positive \emph{magnitude/current law}, not a complex identity. In a general curved metric $\nabla_\mu j^\mu=0$ and the appropriate tetrad and radial measure must be used; \eqref{eq:bridge} is not transplanted there without the geometric factors. If the source is stationary and bound with zero net radial flux, the outgoing division by $p\dot N^{-1}$ does not apply. Current conservation fixes amplitude \emph{relative to a chosen momentum profile}; it supplies neither the source gravity nor a unique value of the original $\alpha_G$.

\section{An exactly solvable Newtonian radial wave}\label{sec:newtonwave}
For comparison with the quantum-wave sector, \emph{assume} an external Newtonian potential $\Phin=-\GNewton M/r$, a positive test fermion mass $m_f$, angular momentum $\ell=0$, and zero nonrelativistic energy. The reduced radial Schr\"odinger equation is
\begin{equation}
 -\frac{\hbar^2}{2m_f}u''+m_f\Phin u=0
 \quad\Longleftrightarrow\quad
 u''+\frac{2\GNewton Mm_f^2}{\hbar^2r}u=0.
 \label{eq:radialSch}
\end{equation}
Define $k=2\GNewton Mm_f^2/\hbar^2>0$ and $z=2\sqrt{kr}$. An exact basis is
\begin{equation}
 \boxed{\displaystyle u(r)=\sqrt r\,[A J_1(z)+B Y_1(z)].}
 \label{eq:Bessel}
\end{equation}
The formally outward-flux Hankel branch is
\begin{equation}
 u_{\rm out}(r)=\mathcal A\sqrt r\,H_1^{(1)}(z),
 \quad R(r)=|\mathcal A|\sqrt r\sqrt{J_1(z)^2+Y_1(z)^2},
 \quad W=\hbar\arg H_1^{(1)}(z).
 \label{eq:Hankel}
\end{equation}
The Bessel Wronskian gives the exact transport identity
\begin{equation}
 \boxed{R(r)^2W'(r)=\frac{\hbar|\mathcal A|^2}{\pi}.}
 \label{eq:exactflux}
\end{equation}
As $z\to\infty$, $R(r)\propto r^{1/4}$ and $W(r)\sim 2m_f\sqrt{2\GNewton Mr}$, reproducing the leading QGD current transport. This is a consistency check \emph{in an assumed Newtonian potential}, not an independent derivation of Newton's law. The Hankel solution is not a globally regular isolated source; the regular $J_1$ standing solution has zero net radial current.

\section{Curved-space Dirac dynamics: the controlled quantum sector}\label{sec:curvedDirac}
To use the classical QGD geometry as a background for a quantum \emph{test} field, make the additional minimal-coupling assumption
\begin{equation}
 (\ii\hbar c\,\Gamma^\mu\nabla_\mu-m_fc^2)\Psi=0,
 \qquad\nabla_\mu=\partial_\mu+\Omega_\mu[g].
 \label{eq:curvedDirac}
\end{equation}
Here the spin connection $\Omega_\mu$ is required. This step is compatible with Chapter 2's constructed Clifford algebra but does not prove the Dirac microscopics generated the metric.

In the Schwarzschild diagonal chart and a specified spherical tetrad/angular spinor basis, define $r_*$ by $\dd r_*/\dd r=1/F$ and use nonzero integer $\kappa$. The separated radial Hamiltonian can be expressed as
\begin{equation}
 \boxed{\begin{pmatrix}
 m_fc^2\sqrt F&-\hbar c\,\dd/\dd r_*+\hbar c\kappa\sqrt F/r\\
 \hbar c\,\dd/\dd r_*+\hbar c\kappa\sqrt F/r&-m_fc^2\sqrt F
 \end{pmatrix}\begin{pmatrix}f\\g\end{pmatrix}
 =E\begin{pmatrix}f\\g\end{pmatrix}.}
 \label{eq:radialDirac}
\end{equation}
The angular spinors and the appropriate spin-connection-dependent prefactor reconstruct the four-component field. Equation \eqref{eq:radialDirac} is a standard \emph{conditional} separated-Dirac result; a fully formalized derivation must independently specify the tetrad, harmonics and gamma-matrix conventions.

\subsection{Exact radial reduction and weak-field coefficients}
Let $\ell_C=\hbar/(m_fc)$, $x=r/\ell_C$, $a=r_s/\ell_C$, $e=E/(m_fc^2)$, $F=1-a/x$, and $\dd/\dd x_*=F\dd/\dd x$. Set $s=\sqrt F$, $B_e=e+s$, $K=\kappa s/x$. The first-order radial system is
\begin{equation}
 f'_* =B_e g-Kf,\qquad g'_* =Kg+(s-e)f.
 \label{eq:radialFirst}
\end{equation}
Eliminate $g$ on a region with $B_e>0$:
\begin{equation}
 f''_{**}-\frac{B'_*}{B_e}f'_*
 +\left[K'_* -\frac{B'_*}{B_e}K-K^2+e^2-F\right]f=0.
 \label{eq:felim}
\end{equation}
The Liouville-normal variable $u_\kappa=\sqrt{F/B_e}\,f$ satisfies
\begin{equation}
 u_\kappa''+\mathcal V_\kappa(x,a,e)u_\kappa=0.
 \label{eq:liouville}
\end{equation}
For $e=1$, $a/x\ll1$ at fixed $x>0$, algebraic expansion gives
\begin{equation}
 \boxed{\displaystyle
 \mathcal V_\kappa=-\frac{\kappa(\kappa+1)}{x^2}
 +\frac ax+\gamma_\kappa\frac a{x^3}+O(a^2),
 \quad\gamma_\kappa=\frac{3-4\kappa^2-\kappa}{4}.}
 \label{eq:Vone}
\end{equation}
For $\kappa=-1$, the first $a/x^3$ term vanishes; the next order is
\begin{equation}
 \boxed{\displaystyle
 \mathcal V_{-1}=\frac ax+\frac{2a^2}{x^2}
 +\frac{a^2}{64x^4}+O(a^3).}
 \label{eq:Vtwo}
\end{equation}
These coefficients are statements about a \emph{particular normalized radial differential equation} and the threshold $e=1$; they are not yet invariant observables or additional field equations for $\sigma_\mu^{(s)}$.

\section{The historical $\hbar^2/r^4$ force: what does and does not follow}\label{sec:quantumforce}
The original QGD program proposed an attractive correction of the form
\begin{equation}
 F_{r,\mathrm{hist}}=-\frac{\GNewton Mm_f}{r^2}
 \left[1+\frac92\left(\frac{\hbar}{m_fc r}\right)^2+\cdots\right].
 \label{eq:histforce}
\end{equation}
This is preserved as a \emph{historical research candidate}, not an established prediction. In the original reciprocal-exponential argument, Taylor-expanding $e^{2\ii m_fcr/\hbar}$ and inverting each nonconstant Taylor term separately is not the reciprocal of the series. Indeed,
\begin{equation}
 \frac{1}{e^{2\ii m_fcr/\hbar}}=e^{-2\ii m_fcr/\hbar},
\end{equation}
and its actual small-$r$ expansion has nonnegative powers of $r$. In addition, the displayed historical cubic energy term yields an outward force, not the more attractive sign in \eqref{eq:histforce}. The regimes $r\ll\ell_C$ (for the exponential Taylor series) and $r\gg\ell_C$ (for a small Compton correction) do not overlap.

One may \emph{formally} rewrite \eqref{eq:liouville} as a one-dimensional Schr\"odinger equation with the radial coefficient $U_\kappa:=-m_fc^2\mathcal V_\kappa/2$. Equations \eqref{eq:Vone}--\eqref{eq:Vtwo} then contain terms
\begin{equation}
 U_\kappa=-\frac{\GNewton Mm_f}{r}
 +\frac{\hbar^2\kappa(\kappa+1)}{2m_fr^2}
 -\gamma_\kappa\frac{\GNewton M\hbar^2}{m_fc^2r^3}
 +O(\GNewton^2).
 \label{eq:radialU}
\end{equation}
The \emph{formal radial-coefficient derivative} is
\begin{equation}
 \delta F_{r,\kappa}^{\rm formal}:=-\frac\dd{\dd r}\left[-\gamma_\kappa
 \frac{\GNewton M\hbar^2}{m_fc^2r^3}\right]
 =-3\gamma_\kappa\frac{\GNewton M\hbar^2}{m_fc^2r^4}.
 \label{eq:formalforce}
\end{equation}
For $\kappa=-1$ it vanishes at first order in $G$, whereas for $\kappa=+1$ it is outward with coefficient $+3/2$. The historical attractive coefficient $-9/2$ would require $\gamma_\kappa=+3/2$, i.e. $4\kappa^2+\kappa+3=0$, whose discriminant is $-47<0$; it does not arise in these channels. At $\kappa=-1$, the first explicitly $\hbar^2$ term at second order produces only the representation-specific
\begin{equation}
 \delta U_{\kappa=-1}^{(G^2\hbar^2)}=-\frac{\hbar^2r_s^2}{128m_fr^4},\qquad
 \delta F_{\kappa=-1}^{\rm formal}=-\frac{\GNewton^2M^2\hbar^2}{8m_fc^4r^5}.
 \label{eq:secondradialforce}
\end{equation}
Neither \eqref{eq:formalforce} nor \eqref{eq:secondradialforce} is asserted to be a physical correction to the gravitational field.

\begin{lemma}[Representation dependence of a Liouville-normal radial coefficient]\label{lem:schwarzian}
If $u''(x)+\mathcal V(x)u(x)=0$ and $x=x(y)$ is a smooth monotone coordinate transformation with $x_y>0$, then the redefined dependent variable $u(x(y))=\sqrt{x_y}\,v(y)$ obeys
\begin{equation}
 v_{yy}+\left[x_y^2\mathcal V(x(y))+\tfrac12\{x,y\}\right]v=0,
 \quad\{x,y\}:=\frac{x_{yyy}}{x_y}-\frac32\left(\frac{x_{yy}}{x_y}\right)^2.
 \label{eq:Schwarzian}
\end{equation}
\end{lemma}
\begin{proof}
Use $\dd/\dd x=x_y^{-1}\dd/\dd y$ and differentiate $u=\sqrt{x_y}\,v$ twice. The coefficient changes by the displayed Schwarzian term. Thus an isolated coefficient of $x^{-3}$ in one chosen radial normal form cannot by itself be identified as a coordinate-independent force on a physical test particle.
\end{proof}

A physical prediction would require a defined observable (e.g. a normalized scattering amplitude, conserved-current response, interferometric phase, or state-dependent acceleration), a specified state and boundary data, and invariance under equivalent tetrad/coordinate choices. A universal quantum \emph{gravitational-field} correction additionally requires a derived quantum-corrected equation for $g[\sigma,T]$ or for the natural invariant itself. The historical standalone $\hbar^2(\nabla\sigma)^2$ action and derivative metric term remain unverified trial operators, not consequences silently imported from Chapters 1--2.

\section{Normalized Dirac sources and where genuine $\hbar$ dependence may enter}\label{sec:source}
The adopted effective action may consistently couple a specified normalized fermionic source to the natural-invariant geometry. In a stationary two-spin-balanced Einstein--Dirac benchmark (in $G=c=\hbar=1$ units) one finds, for an appropriate tetrad,
\begin{equation}
 \Psi_s=\frac{\sqrt{T(r)}}r e^{-\ii\omega t}
 \begin{pmatrix}a(r)\chi_s\\ \ii b(r)(\bm\tau\cdot\widehat{\bm r})\chi_s\end{pmatrix},
 \quad s\in\{\uparrow,\downarrow\},
 \label{eq:sourcewave}
\end{equation}
with real $a,b$ and zero net radial current. This is an \emph{illustrative source benchmark} of adopted Einstein--Dirac dynamics, not a theorem that one coherent four-mode QGD state universally reduces to this form. It demonstrates why the outgoing-flux bridge \eqref{eq:bridge} cannot normalize a stationary self-gravitating source.

To obtain an $\hbar$-dependent \emph{gravitational profile} without postulating a new vacuum gravitational kinetic term, one may calculate the properly normalized quantum matter stress tensor $\langle T_{\mu\nu}\rangle$ and solve its sourced gravitational equations. The profile may depend on $\hbar$ through occupation, gradients or quantum pressure even when the vacuum exterior remains classical Schwarzschild. This is \emph{source-dependent quantum matter gravitating}; it is distinct from a universal quantum-gravity correction to vacuum propagation.


\section{A self-gravitating Dirac-source consistency benchmark}\label{sec:sourcebenchmark}
To distinguish the outgoing test-wave envelope from the macroscopic
\emph{source} envelope, consider a previously solved stationary,
spin-balanced, two-mode Einstein--Dirac configuration.  The statement
below concerns the \emph{adopted effective Einstein--Dirac model}, not
the unrestricted four-mode QGD matter action.
In natural units $G_{\rm N}=c=\hbar=1$, let
\begin{equation}
 \mathrm ds^2=-T(r)^{-2}\mathrm dt^2
    +A_{\rm met}(r)^{-1}\mathrm dr^2+r^2\mathrm d\Omega^2,
 \qquad T>0,\quad A_{\rm met}>0.
 \label{eq:sourceMetric}
\end{equation}
The spin-balanced stationary ansatz may be expressed in a designated
tetrad as
\begin{equation}
 \Psi_s(t,r,\Omega)=\frac{\sqrt{T(r)}}r e^{-\ii\omega t}
 \begin{pmatrix}a_D(r)\chi_s\\[2pt]
 \ii b_D(r)(\bm\tau\cdot\widehat{\bm r})\chi_s\end{pmatrix},
 \qquad s=\uparrow,\downarrow.
 \label{eq:sourceSpinors}
\end{equation}
Here $a_D,b_D$ are the source's two real radial functions, \emph{not}
the two components $A,B$ of the gravitational natural invariant.
The occupied modes are normalized using the proper conserved-current
measure rather than the outward-flux condition:
\begin{equation}
 4\pi\int_0^\infty
 \frac{T(r)}{\sqrt{A_{\rm met}(r)}}
 [a_D(r)^2+b_D(r)^2]\,\mathrm dr=1
 \qquad\text{for each occupied mode}.
 \label{eq:sourceNorm}
\end{equation}
For the chosen local spin frame the upper amplitude in the original QGD
wavefunction may be identified as a \emph{representation-specific}
physical product:
\begin{equation}
 \boxed{\alpha_G^{\rm spin}\psi_0(r)=\sqrt{T(r)}\,a_D(r).}
 \label{eq:sourceUpper}
\end{equation}
This identification leaves the common-rescaling symmetry of
the rescaling discussion of Section~\ref{sec:alpha} untouched: the field equations determine the
physical product, not the arbitrary split into an overall multiplier
and a reference envelope.

For reference, one earlier finite-domain numerical solution used
$m_f=0.3$, total conserved occupation $N=2$, and
$M_{\rm ADM}\simeq0.59893551$ (geometric units).  The source's upper
radial envelope is localized and decays, unlike the outgoing
$r^{1/4}$ wave.  Taking a unit-total-occupation reference gives
$|\alpha_G^{\rm state}|^2=2$, whereas the independent geometric
Compton convention with $m_{\rm ref}=m_f$ gives
$|\alpha_G^{\rm geom}|^2\simeq2.7827$.
The mismatch is \emph{not} a failure of the field equations: it shows
that a reference-amplitude convention and an inverse geometric
strength are different definitions.  Numerical residuals verify
consistency for that finite-domain benchmark but provide no
Lean-certified existence theorem or microscopic promotion proof.

\section{Quantum stress from the Dirac source: a coefficient fixed without new gravity operators}\label{sec:quantumsource}
This section tests a narrowly stated question: can the Dirac matter action already specified in Chapter~1 produce a nonzero coefficient of a quantum source stress \emph{without} inserting a new standalone $\hbar^2(\nabla\sigma)^2$ gravitational operator? The answer is yes in the nonrelativistic positive-energy fixed-spin sector, but this fact by itself does not establish a universal vacuum correction to gravity.

\begin{assumption}[Controlled Dirac-to-Pauli source limit]\label{ass:pauli}
Use the positive-energy sector of the stated Dirac action, a weak static gravitational background supplied by the adopted Chapter~2 effective metric, $|\bm p|\ll m_fc$, no background electromagnetic field, and a locally constant Pauli spinor $\chi$ with $\chi^\dagger\chi=1$. Retain the leading $c^0$ matter Hamiltonian after removing the rest-energy phase. The positive-energy two-component wavefunction is $\varphi=\sqrt{n}\,e^{iS/\hbar}\chi$, where $n>0$ is the \emph{physical} number density of this occupied mode. Spin-texture gradients, antiparticle mixing, relativistic stress corrections, and general many-orbital exchange terms are not included in this single-mode theorem.
\end{assumption}

\begin{proposition}[The Dirac-derived quantum-density gradient coefficient]\label{prop:quantumstress}
Under Assumption~\ref{ass:pauli}, the leading nonrelativistic matter kinetic-energy functional has the exact polar-decomposition identity
\begin{equation}
\boxed{\displaystyle
E_{\rm kin}[n,S]=\int_{\mathbb R^3}\left[
 \frac{n|\bm\nabla S|^2}{2m_f}
 +\frac{\hbar^2}{8m_f}\frac{|\bm\nabla n|^2}{n}\right]\dd^3x.}
\label{eq:gradkin}
\end{equation}
The gradient contribution has functional derivative
\begin{equation}
\boxed{Q[n]=-\frac{\hbar^2}{2m_f}\frac{\nabla^2\sqrt n}{\sqrt n},}
\qquad
\frac{\delta E_{\rm grad}}{\delta n}=Q[n],
\label{eq:Qsource}
\end{equation}
and one symmetric quantum momentum-flux tensor can be chosen as
\begin{equation}
\boxed{\displaystyle
\Pi_{ij}^{Q}=\frac{\hbar^2}{4m_f}
\left(\frac{\partial_i n\,\partial_j n}{n}-\partial_i\partial_jn\right),
\qquad \partial_j\Pi_{ij}^{Q}=n\,\partial_iQ.}
\label{eq:PiQ}
\end{equation}
\end{proposition}
\begin{proof}
For constant $\chi$, $|\bm\nabla(\sqrt n\,e^{iS/\hbar}\chi)|^2=|\bm\nabla\sqrt n|^2+n|\bm\nabla S|^2/\hbar^2$. Multiply by $\hbar^2/(2m_f)$ and integrate to obtain \eqref{eq:gradkin}. Integrate the first variation of $E_{\rm grad}$ by parts with variations vanishing at the boundary to get \eqref{eq:Qsource}. Differentiate \eqref{eq:PiQ} and simplify using the quotient rule and $\sqrt n$; the divergence is $n\bm\nabla Q$. This derivation fixes the $\hbar^2/(8m_f)$ energy coefficient and $\hbar^2/(4m_f)$ stress coefficient. The flux tensor is not unique under divergence-free improvements; its indicated divergence and force are unchanged.
\end{proof}

For this one-mode sector the exact Madelung acceleration equation is
\begin{equation}
 m_fn(\partial_t+\bm v\cdot\bm\nabla)\bm v
 =-m_fn\bm\nabla\Phi-\bm\nabla\cdot\bm\Pi^Q,
 \qquad \bm v=\bm\nabla S/m_f.
 \label{eq:Qhydro}
\end{equation}
The \emph{quantum force on matter} in \eqref{eq:Qhydro} is not itself a new term in the gravitational field equation. The right-hand gravitational acceleration is that of the metric sourced by the matter distribution. Independent occupied orbitals and their stress tensors must be summed to describe a multi-fermion source; a general fermionic mixed state is not characterized solely by one $n$ and $S$.

\begin{proposition}[Gravitational response of the normalized quantum density]\label{prop:quantumsource}
Assume independently the leading nonrelativistic sourced field equation of the adopted Einstein--Hilbert benchmark, $\nabla^2\Phi=4\pi G_{\rm N}\rho$, with $\rho=m_fn+O(c^{-2})$. At this order the dimensionless temporal metric combination in QGD's classical signed-sector representation is $\mathcal U=-2\Phi/c^2$. Thus
\begin{equation}
\boxed{\displaystyle
\nabla^2\mathcal U=-\frac{8\pi G_{\rm N}m_f}{c^2}\,n+O(c^{-4}) .}
\label{eq:Usourced}
\end{equation}
Any $\hbar$-dependence of a \emph{specified and independently solved} quantum state $n(x;\hbar)$ is inherited by $\mathcal U$ through \eqref{eq:Usourced}. No additional coefficient of a QGD vacuum operator has been derived by this argument.
\end{proposition}
\begin{proof}
Substitute $\rho=m_fn$ in the specified Poisson equation and use the stated definition $\mathcal U=-2\Phi/c^2$. The relation is for the temporal combination of the full metric and does not individually fix every signed $\sigma_\mu^{(s)}$ or the spatial metric. At relative order $c^{-2}$, energy density, pressure, gradient stress and gravitational binding require a joint relativistic Einstein--matter calculation; replacing the complete active source by $m_fn+E_{\rm grad}/c^2$ alone would be unjustified.
\end{proof}

\subsection{Exact, normalized source test and the role of Planck's constant}
A clean diagnostic state is the normalized spherically symmetric Gaussian density
\begin{equation}
 n_a(r)=\frac{N}{\pi^{3/2}a^3}e^{-r^2/a^2},
 \qquad \int_{\mathbb R^3}n_a\,\dd^3x=N,
 \qquad M=Nm_f.
 \label{eq:gaussiandensity}
\end{equation}
This is a \emph{test density}, not a claimed exact stationary Dirac-star solution. For it, the quantum potential and gradient energy follow without fitting any coefficient:
\begin{equation}
 Q_a(r)=\frac{\hbar^2}{2m_f}\left(\frac3{a^2}-\frac{r^2}{a^4}\right),
 \qquad E_{\rm grad}(a)=\frac{3N\hbar^2}{4m_fa^2}.
 \label{eq:gaussianQ}
\end{equation}
The unique Newtonian potential vanishing at infinity and solving Poisson's equation for this density is
\begin{equation}
\boxed{\displaystyle
\Phi_a(r)=-\frac{G_{\rm N}M}{r}\,\operatorname{erf}\!\left(\frac ra\right).}
\label{eq:gaussianPhi}
\end{equation}
Its exterior-point radial force on a distinct, slowly moving test mass $m_t$ is
\begin{equation}
\boxed{\displaystyle
 F_r(r)=-\frac{G_{\rm N}Mm_t}{r^2}
 \left[\operatorname{erf}\!\left(\frac ra\right)-\frac{2r}{\sqrt\pi a}e^{-r^2/a^2}\right].}
 \label{eq:gaussianForce}
\end{equation}
This force is a physical gravitational response to the given spatial source distribution; it is distinct from $-\bm\nabla Q$ in \eqref{eq:Qhydro}. The gravitational force contains $\hbar$ only if the matter dynamics independently determine an $\hbar$-dependent width $a$.

For example, in the \emph{specified external} isotropic oscillator model, the Gaussian is an exact nonrelativistic fixed-spin ground-state orbital and its width is $a^2=\hbar/(m_f\omega)$. This yields an exactly calculable $\hbar$-dependent gravitational profile \emph{in that source model}; the confining potential is an extra physical condition and not claimed to follow from Chapter~1 alone.

A separate, strictly variational example avoids assuming a harmonic trap. For two identical spin-$1/2$ fermions in a spin singlet occupying the same trial Gaussian orbital, the expectation of the independent two-particle Newtonian Hamiltonian has no one-particle gravitational self-interaction:
\begin{equation}
 E_{2}(a)=\frac{3\hbar^2}{2m_fa^2}
 -\frac{2G_{\rm N}m_f^2}{\sqrt{2\pi}\,a},
 \qquad a_{\rm var}=\frac{3\sqrt{2\pi}\,\hbar^2}{2G_{\rm N}m_f^3}.
 \label{eq:twofermionvariational}
\end{equation}
The minimum is stable within this one-parameter Gaussian family. The density and classical potential \eqref{eq:gaussianPhi} may be computed for this trial state, but the trial state is \emph{not} an exact coupled Einstein--Dirac stationary solution. For a truly macroscopic many-fermion source, Pauli occupations, exchange and correlations cannot be replaced by putting arbitrarily many fermions into one orbital.

\subsection{Exterior no-go statement for a universal finite-source $r^{-4}$ correction}
\begin{proposition}[Spherical source response outside matter]\label{prop:gaussnoGo}
Let $\rho_\hbar(r)\geq0$ be an integrable, spherically symmetric mass-density profile with fixed total Newtonian mass $M$, and solve $\nabla^2\Phi_\hbar=4\pi G_{\rm N}\rho_\hbar$ with $\Phi_\hbar(\infty)=0$. Then
\begin{equation}
 \boxed{\displaystyle
 -\Phi_\hbar'(r)=-\frac{G_{\rm N}M_\hbar(<r)}{r^2},\quad
 M_\hbar(<r)=4\pi\int_0^r\rho_\hbar(s)s^2\dd s.}
 \label{eq:gaussquantum}
\end{equation}
If $\rho_\hbar$ is supported in $r\leq R$, then for every $r>R$,
$\Phi_\hbar=-G_{\rm N}M/r$ and $-m_t\Phi_\hbar'=-G_{\rm N}Mm_t/r^2$ exactly. A change in the quantum source's \emph{interior profile} at fixed total mass cannot yield an exterior universal correction $\propto G_{\rm N}M\hbar^2/(m_fc^2r^4)$ in this Poisson sector.
\end{proposition}
\begin{proof}
Multiply Poisson's equation by $r^2$, integrate radially using regularity at the center, and apply the boundary value at infinity. For compact support, $M_\hbar(<r)=M$ outside and integrate once more. This does not rule out the state-dependent noncompact Gaussian tail, changes in the total gravitational/ADM mass, relativistic source corrections, nonspherical multipoles, or an independently derived modification of the gravitational response law. For the Gaussian, the bracket in \eqref{eq:gaussianForce} differs from one exponentially at large fixed $r/a$; it is not a universal power series in $\hbar^2/r^2$.
\end{proof}

\paragraph{Status of the sought quantum response coefficient.}
The coefficient $\hbar^2/(8m_f)$ multiplying the density gradient is \emph{derived} within the stated Dirac-to-Pauli source sector, and $\Pi^Q$ follows by variation. It is not a coefficient of a gravitational natural-invariant propagation operator. The original proposed $\hbar^2(\nabla\sigma)^2$ and $\lambda\mathcal Q_{ab}$ require their own microscopic matching calculation and independent variation before they can alter the physical vacuum gravitational response. In particular, nothing in the calculation fixes the historical universal attractive coefficient $9/2$; it cannot be imported by Newtonian normalization.


\section{Action-level handoff to Chapter 4}\label{sec:handoff4}
The normalization theorem of Sections~\ref{sec:static}--\ref{sec:series}
determines the exterior constant $C_g$ and fixes the classical static
Killing-energy function; it does not by itself supply the unrestricted
gravitational equations for the independent natural-invariant fields.
To establish the precise logical handoff, define in the \emph{adopted}
metric-only $\lambda=0$ fixed-frame effective action
\begin{equation}
 \Gamma[g,\Psi]=S_{\rm EH}[g]+S_D[g,\Psi],\qquad
 \mathcal E_g^{\mu\nu}:=\frac{1}{\sqrt{-g}}
 \frac{\delta\Gamma}{\delta g_{\mu\nu}}.
 \label{eq:handoffmetric}
\end{equation}
For the signed specialization
$g_{\mu\nu}=\bar g_{\mu\nu}+\sum_s s_s\sigma^{(s)}_\mu\sigma^{(s)}_\nu$,
independent variations at fixed frame and fixed other sectors yield
\begin{equation}
 \boxed{\displaystyle
 \mathcal E_{(s)}^\rho:=\frac{1}{\sqrt{-g}}
 \frac{\delta\Gamma}{\delta\sigma_\rho^{(s)}}
 =2s_s\mathcal E_g^{\rho\nu}\sigma_\nu^{(s)}.}
 \label{eq:handoffprojected}
\end{equation}
This chain-rule equation, which Chapter~4 proves in detail, is not an
additional $\hbar$-dependent wave equation. When the derivative of the
metric map has rank ten for admissible signed-sector variations, the
complete collection of equations~\eqref{eq:handoffprojected} is
\emph{equivalent to the adopted effective metric equations}. With
insufficient rank it is only a projection. An independent $\sigma$
kinetic operator or a nonzero derivative-metric term is not inferred
from the notation $\sigma_\mu$ or from the success of the Schwarzschild
solution.

The leading Dirac-to-Pauli gradient stress in
Proposition~\ref{prop:quantumstress} contributes to $T_{\mu\nu}$ in a
controlled sourced approximation; it does \emph{not} create an
independent $\sigma$ kinetic term in~\eqref{eq:handoffprojected}.
The Poisson exterior result of Proposition~\ref{prop:gaussnoGo}
precludes a universal power-law $\hbar^2/r^4$ force from changing
only the interior profile of a compact spherical source at fixed mass.
A genuine QGD-specific exterior quantum-force claim therefore requires
an independently varied quantum effective action or new source-to-natural-
invariant constitutive dynamics, along with a physical observable.

\section{What Chapter 3 proves, and formal verification gates}\label{sec:proofledger}
\begin{center}
\begin{tabular}{>{\raggedright\arraybackslash}p{0.54\textwidth}>{\raggedright\arraybackslash}p{0.34\textwidth}}
\toprule
Statement & Logical status\\\midrule
Microscopic Dirac-to-macroscopic four-mode structure & Conditional local/WKB construction\\
Promotion to the independent $\sigma_{\mu}^{(s)}$ fields & QGD postulate\\
Complete signed/tetrad/derivative metric & Proposed definition; $\mathcal Q$ unverified\\
Spherical vacuum $A^2=B^2=C_g/r$ & Derived in explicitly adopted EH sector\\
Newtonian mass calibration $C_g=2GM/c^2$ & Permissible one-parameter match\\
Exact static potential and all Taylor coefficients & Derived classical benchmark\\
Geometric $\alpha_G^{\rm geom}$ & Explicit dimensionless convention\\
Identification with original spinor $\alpha_G$ & Open, not proven\\
Flux bridge for outgoing Dirac WKB waves & Derived conditional transport identity\\
Newtonian Hankel wave & Exact for prescribed Newtonian potential\\
Curved-space separated Dirac coefficients & Conditional; algebraically auditable\\
Dirac-derived quantum gradient source coefficient $\hbar^2/(8m_f)$ & Derived under fixed-spin Pauli hypotheses\\
Spherical quantum-source gravity and Gaussian test & Conditional Poisson source theorem; $\hbar$ enters through state\\
Universal exterior $r^{-4}$ from compact spherical source alone & Excluded at fixed mass in Poisson sector\\
Universal historical attractive $9/2$ quantum force & Not established\\
Metric-only $\sigma$ action chain rule & Derived; Chapter 4 full-rank conditions required\\
Lean 4 kernel verification & Pending, not claimed\\
\bottomrule
\end{tabular}
\end{center}

\subsection*{Formalization obligations}
Before claiming Lean verification, formal proofs must cover: (i) positivity and domain of $F,D,r$; (ii) exact four-dimensional Ricci and boundary identities; (iii) both variational derivatives with the boundary functional; (iv) independent ODE integration and asymptotic hypotheses; (v) Newtonian correspondence assumptions; (vi) exact static energy and analytic series with an explicit tail bound; (vii) normalization-rescaling invariance and $\hbar$ cancellation; (viii) the WKB current identity, with Pauli/gamma matrices; (ix) the stated Newtonian Bessel differential equation and Wronskian; (x) the conditional separated Dirac system and Liouville transformation; and (xi) the Schwarzian representation-dependence lemma; (xii) Dirac-to-Pauli sector hypotheses, polar kinetic decomposition and quantum-stress divergence; and (xiii) the Poisson spherical-source exterior theorem; (xiv) the metric-to-natural-invariant functional chain rule and admissible rank-ten variations. The computational checks accompanying this chapter validate selected algebra, \emph{not} the omitted analytic, boundary-value, microscopic or tensor-calculus proof obligations.

\section{Conclusion}
The chapter's established conditional result is a \emph{Newton-calibrated classical} gravitational natural-invariant exterior with the exact static Killing-energy potential \eqref{eq:PhiExact} and the rigorously controlled expansion \eqref{eq:fourterms}. It also exhibits a calculable quantum Dirac-wave sector with explicit $\hbar$ and a correct radial-current bridge. Neither the historical inverse-exponential derivation nor the separated Dirac radial coefficient produces a verified universal $\hbar^2/r^4$ gravitational force. A Dirac-derived quantum-gradient stress with the fixed coefficient $\hbar^2/(8m_f)$ has now been established in a stated positive-energy fixed-spin source limit, and its Poisson-sourced gravitational response and exterior restrictions were checked explicitly. This does not prove a universal quantum correction to gravitational propagation. A future universal quantum-force claim requires an independently derived source-to-$\sigma$ response or gravitational kinetic coefficient, together with a physical invariant observable and successful limits.




\section*{Q061 final interpretation of normalization and quantum order}
\addcontentsline{toc}{section}{Q061 final normalization and quantum-order audit}
\noindent Newtonian matching fixes the \emph{single} measured gravitational coupling $G_N$ in the adopted Einstein--Hilbert metric-mediated action. It does not require the independent microscopic existence postulate (Axiom II) to have been proved. The exact static Killing-energy potential and its $r^{-2}$ and $r^{-3}$ Taylor terms on Schwarzschild are classical nonlinearities, even when they resemble proposed high-order quantum corrections. A correction is identified as quantum here only by its provenance in the consistently renormalized state/loop response (or explicitly stated quantum source), not merely by its radial order.

\noindent In the present effective-action organization the local $R^2,C^2,E_4$ couplings are renormalized matching coefficients; QN004 supplies a calculated one-loop contribution to the momentum-dependent curvature kernels. The same contribution must not be reinserted into the natural-invariant master-metric slot as an additional $\mathcal Q_{ab}$ interaction. The Chapter~7 far-field massless-Dirac Boulware example illustrates the distinction: it provides conditional order-$\hbar$ changes to the physical metric and its selected natural-invariant representation, rather than a universal new term in the Newtonian potential. The normalizations and original wavefunction results in this chapter are unchanged.


\appendix
\chapter{Proof details and reproducibility ledger}\label{app:proof}
\section{Independent variation without fixing the natural-invariant branch}\label{app:EL}
For clarity, define $F=1-A^2$, $D=1+B^2-A^2$ and
$L_{\rm red}=[F+rF']/\sqrt D+\sqrt D$.
The Euler--Lagrange derivatives are defined by
\begin{equation}
 \mathcal E_A:=\frac{\partial L_{\rm red}}{\partial A}
 -\frac{\mathrm d}{\mathrm dr}
  \frac{\partial L_{\rm red}}{\partial A'},
 \qquad
 \mathcal E_B:=\frac{\partial L_{\rm red}}{\partial B}
 -\frac{\mathrm d}{\mathrm dr}
  \frac{\partial L_{\rm red}}{\partial B'}.
\end{equation}
Because $L_{\rm red}$ depends on $B$ algebraically, the second
expression is an ordinary partial derivative.  Substituting
$F'=-2AA'$ and $D'=2BB'-2AA'$ and simplifying yields
\begin{equation}
 \mathcal E_A=-\frac{A(2rBB'+B^2)}{D^{3/2}},
 \qquad
 \mathcal E_B=\frac{B(B^2+2rAA')}{D^{3/2}}.
\end{equation}
These reproduce Eq.~\eqref{eq:EL} and show explicitly that both
components are varied before the equality $A^2=B^2$ is obtained.
The construction is valid only under the declared differentiability,
nondegeneracy, boundary and action hypotheses.

\section{Convergent coefficient and error proof}\label{app:series}
For $n\ge1$ let $a_n={2n\choose n}/[4^n(2n-1)]$.
The elementary identity
\begin{equation}
 \frac{a_{n+1}}{a_n}=\frac{2n-1}{2n+2}
\end{equation}
shows $a_n$ is positive and decreasing.  For $0<x<1$,
$|\sum_{n=5}^\infty a_nx^n|\le a_5x^5/(1-x)$ with
$a_5=7/256$.  This proves the exact remainder inequality in
Theorem~\ref{thm:series} and makes the convergence domain and
power-counting assumptions explicit.

\section{Formal-verification boundary}\label{app:lean}
A Lean kernel proof has \emph{not} been executed.  The independent
symbolic tests check mathematical identities, not the derivation of
the microscopic effective action, the admissibility of the QGD
promotion postulate, or the existence of the numerical boundary-value
states.  A full kernel proof would need, in logical order:
\begin{enumerate}
 \item Positivity, differentiability and boundary assumptions for the
 spherical metric functions, radial integrals and variations;
 \item A formal Levi--Civita connection, Ricci contraction and exact
 boundary identity of Lemma~\ref{lem:EH};
 \item Euler--Lagrange derivatives of both $A$ and $B$ and the
 nondegenerate $C_g/r$ exterior integration;
 \item Axiomatization of Newtonian correspondence, static Killing
 energy and exact Taylor coefficient identities with a tail bound;
 \item Dirac Pauli-matrix current identities, their state/transport
 qualifications and the reference-normalization symmetry;
 \item The Bessel ODE and Wronskian for the prescribed Newtonian
 radial equation;
 \item Spherical spin-connection separation, Liouville transformation
 and weak-field Dirac coefficient expansion for specified parameters;
 \item The Schwarzian field-redefinition lemma and the operational
 distinction between one-dimensional coefficients and invariant
 gravitational observables.
\end{enumerate}
None of these obligations licenses an inference from a classical test
Dirac radial correction to a universal vacuum quantum-gravity force.


\clearpage
\end{document}
